%! TEX root = ../bcd-19-20.tex

\chapter{Extreme cu legături}

\section{Generalități}

În cazul în care domeniul de definiție este specificat cu o ecuație sau este
dat de un produs de intervale, se aplică o metodă specifică, numită, în general,
\emph{metoda multiplicatorilor lui Lagrange}.

Astfel, fie funcția $ f : D \seq \RR^2 \to \RR $ ca mai sus, căreia vrem să îi determinăm
extremele și presupunem că vrem să o facem numai într-un domeniu dat de:
\[
  D = \{ (x, y) \in \RR^2 \mid x^2 + y^2 = 4 \}.
\]

Atunci, ecuația $ x^2 + y^2 = 4 $ se numește \emph{legătură} și o scriem sub
forma unei funcții $ g(x, y) = x^2 + y^2 - 4 $, pentru ca legătura să devină
$ g(x, y) = 0 $.

Cu aceasta, metoda multiplicatorilor lui Lagrange înseamnă să alcătuim
\emph{funcția lui Lagrange}\footnotemark:
\[
  F(x, y) = f(x, y) - \lambda g(x, y), \quad \lambda \in \RR.
\]
\index{funcții!Lagrange}
\index{funcții!multiplicator Lagrange}
\footnotetext{Detalii și explicații geometrice sînt date, de exemplu, %
    foarte clar în această %
\href{https://www.youtube.com/watch?v=CQnZy4n85fE}{lecție video}.}

Cu aceasta, singura modificare care trebuie aplicată metodei de mai sus este
că sistemul pentru găsirea punctelor critice devine:
\[
  \begin{cases}
    f_x &= 0 \\
    f_y &= 0 \\
    g &= 0
  \end{cases} \Longrightarrow (x, y, \lambda).
\]
Pentru punctele critice, ne putem dispensa de $ \lambda $, iar rezolvarea funcționează
ca mai devreme.

\begin{remark}\label{rk:leg-simpla}
  Nu întotdeauna este nevoie să ne complicăm cu funcția lui Lagrange.
  Dacă, de exemplu, legătura era $ x + y = 3 $, putem să o scriem ca $ y = 3 - x $,
  iar funcția inițială devine o funcție de o variabilă, $ f(x, 3 - x) $,
  căreia îi studiem extremele ca în liceu.
\end{remark}

\begin{remark}\label{rk:mm-leg}
  Dacă legăturile sînt multiple, de exemplu, $ g_1, g_2, g_3, \dots, g_k $,
  funcția lui Lagrange corespunzătoare este:
  \[
    F(x, y) = f(x, y) - \sum \lambda_i g_i(x, y),
  \]
  iar sistemul de rezolvat este dat de anularea derivatelor parțiale și
  a tuturor legăturilor.
\end{remark}

%%%%%%%%%%%%%%%%%%%%%%%%%%%%%%%%%%%%%%%%%%%%%%%%%%%%%%%%%%%%%%%%%%%%%% 

\subsection{Cazul compact}

Dacă domeniul de definiție este un \emph{spațiu compact} --- ceea ce, în esență,
pentru uzul acestui seminar înseamnă un produs de intervale (semi)închise
sau legături date de inegalități ---, problema se studiază în două etape:
\begin{itemize}
\item \textbf{În interiorul domeniului}, caz în care legătura este inexistentă;
\item \textbf{Pe frontieră}, caz în care legătura este dată de egalitate.
\end{itemize}

De exemplu, pentru domeniile:
\[
  D_1 = [3, 4] \times [1, 5], \quad D_2 = \{ (x, y) \in \RR^2 \mid 3x + 2y \leq 2 \}
\]
\begin{itemize}
\item Interiorul înseamnă:
  \begin{itemize}
  \item pentru $ D_1 $, $ \dr{Int}D_1 = (3, 4) \times (1, 5) $, caz în care avem
    doar de verificat că punctele critice se găsesc în intervalele date;
  \item pentru $ D_2 $, avem legătura $ 3x + 2y < 2 $, caz în care, din nou,
    verificăm dacă punctele critice satisfac inegalitatea \emph{strictă}.
  \end{itemize}
\item Frontierele înseamnă:
  \begin{itemize}
  \item pentru $ D_1 $, cazurile separate $ x = 3, y \in [1, 5] $, apoi
    $ x = 4, y \in [1, 5] $ și invers;
  \item pentru $ D_2 $, legătura devine $ 3x + 2y = 2 $, pe care o putem rezolva
    cu Lagrange sau cu metoda simplificată din observația~\ref{rk:leg-simpla}.
  \end{itemize}
\end{itemize}

%%%%%%%%%%%%%%%%%%%%%%%%%%%%%%%%%%%%%%%%%%%%%%%%%%%%%%%%%%%%%%%%%%%%%%

\section{Exerciții}

\indent\indent 1. Să se determine valorile extreme ale funcțiilor $ f $,
cu legătura $ g $ în cazurile:
\begin{enumerate}[(a)]
\item $ f(x, y) = x^2 + y^2 - 6x - 2y + 1, g(x, y) = x + y - 2 $;
\item $ f(x, y) = x^2 + y^2 - 2y + 1, g(x, y) = x^2 + y^2 - 2x - 2y + 1 $;
\item $ f(x, y) = 3x + 4y, g(x, y) = x^2 + y^2 + 25 $
\end{enumerate}

\vspace{1cm}

2. Să se determine valorile extreme pentru funcțiile $ f $, definite pe
mulțimea $ D $:
\begin{enumerate}[(a)]
\item $ f(x, y) = 5x^2 + 4xy + 8y^2, D = \{(x, y) \in \RR^2 \mid x^2 + y^2 \leq 1 \} $;
\item $ f(x, y) = xy^2(x + y - 2), D = \{(x, y) \in \RR^2 \mid x + y \leq 3, x, y \geq 0 \} $;
\item $ f(x, y) = \dfrac{x + y}{1 + xy}, D = [0, 1] \times [0, 1] $;
\item $ f(x, y, z) = x^2 + y^2 + xz - yz, D = \{(x, y, z \in \RR^3 \mid %
  x + y + z = 1, x, y, z \geq 0 \} $.
\end{enumerate}

%%% Local Variables:
%%% mode: latex
%%% TeX-master: "../bcd-20-21"
%%% End:
